% !PRINTED 6
\weberpage{6}

\gesperrt{5. Gruppen, die diese Eigenschaft haben, heissen commutative Gruppen oder auch Abel'sche Gruppen.}

Wenn sich die Elemente zweier Gruppen
\[
a, b, c, d \ldots
\]
und
\[
a', b', c', d' \ldots
\]
in der Weise gegenseitig eindeutig entsprechen, dass immer, wenn $ab = c$ ist, auch $a'b' = c'$ wird, so heissen die Gruppen \gesperrt{isomorph}, und es gilt der evidente Satz, dass zwei mit einer dritten isomorphe Gruppen unter einander isomorph sind. Man kann hiernach alle unter einander isomorphe Gruppen zu einer Classe von Gruppen zusammenfassen, die selbst wieder eine Gruppe ist, deren Elemente die Gattungsbegriffe sind, die man erhält, wenn man die entsprechenden Elemente der einzelnen isomorphen Gruppen zu einem Allgemeinbegriff zusammenfasst. Die einzelnen unter einander isomorphen Gruppen sind dann als verschiedene Repräsentanten eines Gattungsbegriffes aufzufassen.

Die Eigenschaften der Gruppen, die in Betracht kommen können, sind von verschiedener Art. Sie können nämlich entweder den besonderen Gruppen anhaften und aus der Natur der Elemente abgeleitet sein, aus denen die Gruppe besteht, oder auch aus der Natur des Compositionsgesetzes. Oder sie können den Gruppen als solchen anhaften und müssen sich dann lediglich aus der Definition des Gruppenbegriffes ableiten lassen. Die letzteren Eigenschaften kommen allen isomorphen Gruppen gemeinsam zu und können als invariante Eigenschaften der Gruppe bezeichnet werden. Wenn ein Vergleich gestattet ist, so könnte man an den Unterschied zwischen den metrischen und projectiven Eigenschaften in der Geometrie erinnern.

Zu der ersten Art der Eigenschaften, die aus der besonderen Natur der Elemente abgeleitet werden, gehören z.\,B. bei den Permutationsgruppen die Eigenschaften der Transitivität und Intransitivität, der Primitivität und Imprimitivität; zu den invarianten Eigenschaften gehören die Vertauschbarkeit oder Nichtvertauschbarkeit, der Grad, die Divisoren und ihr Index, die Normaltheiler. Mit diesen invarianten Eigenschaften, bei denen es gleichgültig ist, aus welchem Repräsentanten einer Classe isomorpher Gruppen sie abgeleitet sind, haben wir uns zunächst zu beschäftigen.
